% GENERATED FILE. Edit canonical lessons, then run npm run generate:books.
% canonical-source-hash: fnv1a64:ab90c4913eea21f3
% canonical-lessons: BN-C423-thonga, BN-C423-morok, BN-C423-keribyag, BN-C423-saij, BN-C423-gochhano, BN-R423-review-416-423

\chapter{Bags, Wrappers and Sizes}
\label{ch:bn-bags-wrappers-and-sizes}
\hlchaptermodality{423}

\begin{chapteropening}
\textbf{By the end of this chapter:} \emph{I can talk about bags and wrappers, ask for my size and say I was palmed off.}

Complete the last lesson of chapter 423: I can talk about bags and wrappers, ask for my size and say I was palmed off.
\end{chapteropening}

\section[\texorpdfstring{\bn{ঠোঙা} (ṭhongā)}{ঠোঙা (ṭhongā)}]{\bn{ঠোঙা} (ṭhongā) — a paper bag; a paper cone}

\label{lesson:BN-C423-thonga}



\begin{quote}
Before the new one: say the Bengali for a customer (formal), then the Bengali for to test (by trying).
\end{quote}

\subsection*{You'll want to know: \bn{ঠোঙা}}
\textbf{\bn{ঠোঙা}} — \emph{ṭhongā} — \textquotedblleft{}a paper bag; a paper cone\textquotedblright{}.

\begin{quote}
\bn{মুদি} \bn{একটা} \bn{ঠোঙায়} \bn{মুড়ি} \bn{দিলেন।}
\end{quote}

\begin{quote}
\emph{mudi ekṭā ṭhongāy muṛi dilen.}
\end{quote}

\begin{quote}
The grocer gave me puffed rice in a paper bag.
\end{quote}

Folded and pasted from old newspaper, the \bn{ঠোঙা} holds roasted nuts, puffed rice and sweets.

\begin{cousinweb}
\textbf{\bn{ঠোঙা}} (\emph{ṭhongā}) is the paper bag of every Bengal shop; its origin is not settled.
\end{cousinweb}

\begin{grammarlens}[title={What you've built}]
One more word for getting around.
\end{grammarlens}

\subsection*{Your turn}
\begin{itemize}
\raggedright
  \item \emph{Say it:} \emph{ṭhongā}
  \item \emph{Say it:} \emph{ṭhongā}, once more
  \item \emph{Say it:} \emph{mudi ekṭā ṭhongāy muṛi dilen.} — The grocer gave me puffed rice in a paper bag.
  \item \emph{Recall:} say the Bengali for a customer (formal), then the Bengali for to test (by trying), then say \emph{ṭhongā} again
\end{itemize}

\begin{tcolorbox}[breakable,colback=teal!4,colframe=teal!35!black,title={Before you move on}]
What does \bn{ঠোঙা} mean? (\textquotedblleft{}A paper bag; a paper cone\textquotedblright{}.) Say it once more.
\end{tcolorbox}
\section[\texorpdfstring{\bn{মোড়ক} (moṛôk)}{মোড়ক (moṛôk)}]{\bn{মোড়ক} (moṛôk) — a wrapper; packaging}

\label{lesson:BN-C423-morok}



\begin{quote}
Before the new one: say the Bengali for to test (by trying), then the Bengali for a paper bag.
\end{quote}

\subsection*{You'll want to know: \bn{মোড়ক}}
\textbf{\bn{মোড়ক}} — \emph{moṛôk} — \textquotedblleft{}a wrapper; packaging\textquotedblright{}.

\begin{quote}
\bn{আমি} \bn{উপহারের} \bn{মোড়ক} \bn{খুললাম।}
\end{quote}

\begin{quote}
\emph{āmi uphārer moṛôk khullām.}
\end{quote}

\begin{quote}
I opened the wrapping of the present.
\end{quote}

\begin{cousinweb}
\textbf{\bn{মোড়ক}} (\emph{moṛôk}) is built on \textbf{\bn{মোড়া}}, \textquotedblleft{}to wrap\textquotedblright{}, which CDIAL sets beside \emph{mōṭati}, \textquotedblleft{}twists\textquotedblright{} (CDIAL 10186), with the noun ending -\bn{ক}: the thing that wraps.
\end{cousinweb}

\begin{grammarlens}[title={What you've built}]
One more word for getting around.
\end{grammarlens}

\subsection*{Your turn}
\begin{itemize}
\raggedright
  \item \emph{Say it:} \emph{moṛôk}
  \item \emph{Say it:} \emph{moṛôk}, once more
  \item \emph{Say it:} \emph{āmi uphārer moṛôk khullām.} — I opened the wrapping of the present.
  \item \emph{Recall:} say the Bengali for to test (by trying), then the Bengali for a paper bag, then say \emph{moṛôk} again
\end{itemize}

\begin{tcolorbox}[breakable,colback=teal!4,colframe=teal!35!black,title={Before you move on}]
What does \bn{মোড়ক} mean? (\textquotedblleft{}A wrapper; packaging\textquotedblright{}.) Say it once more.
\end{tcolorbox}
\section[\texorpdfstring{\bn{ক্যারিব্যাগ} (kæribyāg)}{ক্যারিব্যাগ (kæribyāg)}]{\bn{ক্যারিব্যাগ} (kæribyāg) — a carrier bag; a plastic bag}

\label{lesson:BN-C423-keribyag}



\begin{quote}
Before the new one: say the Bengali for a paper bag, then the Bengali for a wrapper.
\end{quote}

\subsection*{You'll want to know: \bn{ক্যারিব্যাগ}}
\textbf{\bn{ক্যারিব্যাগ}} — \emph{kæribyāg} — \textquotedblleft{}a carrier bag; a plastic bag\textquotedblright{}.

\begin{quote}
\bn{দোকানদার} \bn{একটা} \bn{ক্যারিব্যাগে} \bn{জিনিসগুলো} \bn{দিলেন।}
\end{quote}

\begin{quote}
\emph{dokāndār ekṭā kæribyāge jinisgulo dilen.}
\end{quote}

\begin{quote}
The shopkeeper gave me the things in a carrier bag.
\end{quote}

\begin{cousinweb}
\textbf{\bn{ক্যারিব্যাগ}} (\emph{kæribyāg}) is borrowed from English \emph{carry bag}, the Indian English name for a shop's plastic bag, written as one word; \bn{ব্যাগ} inside it is already taught.
\end{cousinweb}

\begin{grammarlens}[title={What you've built}]
One more word for getting around.
\end{grammarlens}

\subsection*{Your turn}
\begin{itemize}
\raggedright
  \item \emph{Say it:} \emph{kæribyāg}
  \item \emph{Say it:} \emph{kæribyāg}, once more
  \item \emph{Say it:} \emph{dokāndār ekṭā kæribyāge jinisgulo dilen.} — The shopkeeper gave me the things in a carrier bag.
  \item \emph{Recall:} say the Bengali for a paper bag, then the Bengali for a wrapper, then say \emph{kæribyāg} again
\end{itemize}

\begin{tcolorbox}[breakable,colback=teal!4,colframe=teal!35!black,title={Before you move on}]
What does \bn{ক্যারিব্যাগ} mean? (\textquotedblleft{}A carrier bag; a plastic bag\textquotedblright{}.) Say it once more.
\end{tcolorbox}
\section[\texorpdfstring{\bn{সাইজ} (sāij)}{সাইজ (sāij)}]{\bn{সাইজ} (sāij) — a size (of clothes or shoes)}

\label{lesson:BN-C423-saij}



\begin{quote}
Before the new one: say the Bengali for a wrapper, then the Bengali for a carrier bag.
\end{quote}

\subsection*{You'll want to know: \bn{সাইজ}}
\textbf{\bn{সাইজ}} — \emph{sāij} — \textquotedblleft{}a size (of clothes or shoes)\textquotedblright{}.

\begin{quote}
\bn{এই} \bn{জুতোটা} \bn{আমার} \bn{সাইজের} \bn{নয়।}
\end{quote}

\begin{quote}
\emph{ei jutoṭā āmār sāijer nôy.}
\end{quote}

\begin{quote}
These shoes are not my size.
\end{quote}

\begin{cousinweb}
\textbf{\bn{সাইজ}} (\emph{sāij}) is borrowed from English \emph{size}; the English z is said with \bn{জ}.
\end{cousinweb}

\begin{grammarlens}[title={What you've built}]
One more word for getting around.
\end{grammarlens}

\subsection*{Your turn}
\begin{itemize}
\raggedright
  \item \emph{Say it:} \emph{sāij}
  \item \emph{Say it:} \emph{sāij}, once more
  \item \emph{Say it:} \emph{ei jutoṭā āmār sāijer nôy.} — These shoes are not my size.
  \item \emph{Recall:} say the Bengali for a wrapper, then the Bengali for a carrier bag, then say \emph{sāij} again
\end{itemize}

\begin{tcolorbox}[breakable,colback=teal!4,colframe=teal!35!black,title={Before you move on}]
What does \bn{সাইজ} mean? (\textquotedblleft{}A size (of clothes or shoes)\textquotedblright{}.) Say it once more.
\end{tcolorbox}
\section[\texorpdfstring{\bn{গছানো} (gôchhāno)}{গছানো (gôchhāno)}]{\bn{গছানো} (gôchhāno) — to palm off; to foist (something on someone)}

\label{lesson:BN-C423-gochhano}



\begin{quote}
Before the new one: say the Bengali for a carrier bag, then the Bengali for a size (of clothes or shoes).
\end{quote}

\subsection*{You'll want to know: \bn{গছানো}}
\textbf{\bn{গছানো}} — \emph{gôchhāno} — \textquotedblleft{}to palm off; to foist (something on someone)\textquotedblright{}.

\begin{quote}
\bn{দোকানদার} \bn{আমাকে} \bn{একটা} \bn{পচা} \bn{আম} \bn{গছাল।}
\end{quote}

\begin{quote}
\emph{dokāndār āmāke ekṭā pôchā ām gôchhālô.}
\end{quote}

\begin{quote}
The shopkeeper palmed off a rotten mango on me.
\end{quote}

Past (he/she …): \bn{তিনি} \bn{পুরনো} \bn{মাল} \bn{গছালেন}, \emph{tini purôno māl gôchhālen}.

The person who is landed with it takes -\bn{কে}: \bn{আমাকে} \bn{গছাল}, foisted it on me.

\begin{cousinweb}
\textbf{\bn{গছানো}} (\emph{gôchhāno}) probably belongs with Sanskrit \emph{gácchati}, \textquotedblleft{}goes, falls to the share of\textquotedblright{} (CDIAL 3955), under which CDIAL lists Bengali \emph{gachāna}, \textquotedblleft{}to take a deposit\textquotedblright{}; to palm off is to make a thing fall to someone else's share.
\end{cousinweb}

\begin{grammarlens}[title={What you've built}]
One more word for getting around.
\end{grammarlens}

\subsection*{Your turn}
\begin{itemize}
\raggedright
  \item \emph{Say it:} \emph{gôchhāno}
  \item \emph{Say it:} \emph{gôchhāno}, once more
  \item \emph{Say it:} \emph{dokāndār āmāke ekṭā pôchā ām gôchhālô.} — The shopkeeper palmed off a rotten mango on me.
  \item \emph{Recall:} say the Bengali for a carrier bag, then the Bengali for a size (of clothes or shoes), then say \emph{gôchhāno} again
\end{itemize}

\begin{tcolorbox}[breakable,colback=teal!4,colframe=teal!35!black,title={Before you move on}]
What does \bn{গছানো} mean? (\textquotedblleft{}To palm off; to foist (something on someone)\textquotedblright{}.) Say it once more.
\end{tcolorbox}
\section[(dialogue)]{Looking back: from Debts and Instalments to Bags, Wrappers and Sizes}

\label{lesson:BN-R423-review-416-423}



\begin{quote}
Say the Bengali for an instalment and a debt (owed).
\end{quote}

\subsection*{Your turn}
Say these aloud:

\begin{itemize}
\raggedright
  \item an instalment, a debt (owed), money owed to one, to pay off, to settle (a bill)
  \item a banknote, a cash machine, a (bank) balance, to transfer (money), a contribution (to a fund)
  \item profit (in trade), a duty (tax), land tax, a moneylender, the rest (remaining)
  \item want (lack), waste (of something useful), stingy, a pledge (for a loan), one's means
  \item a sale (selling), a hawker, a broker, a village market, to be sold
  \item goods, a warehouse, an order (for goods or food), wholesale, to stock up
  \item a pair of scales, pure (unmixed), a shopping list, a customer (formal), to test (by trying)
  \item a paper bag, a wrapper, a carrier bag, a size (of clothes or shoes), to palm off
\end{itemize}

\begin{tcolorbox}[breakable,colback=teal!4,colframe=teal!35!black,title={Before you move on}]
An instalment? (\textbf{\bn{কিস্তি}}, \emph{kisti}.) A debt (owed)? (\textbf{\bn{দেনা}}, \emph{denā}.) Money owed to one? (\textbf{\bn{পাওনা}}, \emph{pāonā}.) To pay off? (\textbf{\bn{শোধ} \bn{করা}}, \emph{shodh kôrā}.) To settle (a bill)? (\textbf{\bn{মেটানো}}, \emph{meṭāno}.) A banknote? (\textbf{\bn{নোট}}, \emph{noṭ}.) A cash machine? (\textbf{\bn{এটিএম}}, \emph{eṭiem}.) A (bank) balance? (\textbf{\bn{ব্যালান্স}}, \emph{bælāns}.) To transfer (money)? (\textbf{\bn{ট্রান্সফার} \bn{করা}}, \emph{ṭrānsphār kôrā}.) A contribution (to a fund)? (\textbf{\bn{চাঁদা}}, \emph{chā̃dā}.) Profit (in trade)? (\textbf{\bn{মুনাফা}}, \emph{munāphā}.) A duty (tax)? (\textbf{\bn{শুল্ক}}, \emph{shulkô}.) Land tax? (\textbf{\bn{খাজনা}}, \emph{khājnā}.) A moneylender? (\textbf{\bn{মহাজন}}, \emph{môhājôn}.) The rest (remaining)? (\textbf{\bn{বাকি}}, \emph{bāki}.) Want (lack)? (\textbf{\bn{অভাব}}, \emph{ôbhāb}.) Waste (of something useful)? (\textbf{\bn{অপচয়}}, \emph{ôpôchôy}.) Stingy? (\textbf{\bn{কিপটে}}, \emph{kipṭe}.) A pledge (for a loan)? (\textbf{\bn{বন্ধক}}, \emph{bôndhôk}.) One's means? (\textbf{\bn{সম্বল}}, \emph{sômbôl}.) A sale (selling)? (\textbf{\bn{বিক্রি}}, \emph{bikri}.) A hawker? (\textbf{\bn{ফেরিওয়ালা}}, \emph{pherioyālā}.) A broker? (\textbf{\bn{দালাল}}, \emph{dālāl}.) A village market? (\textbf{\bn{হাট}}, \emph{hāṭ}.) To be sold? (\textbf{\bn{বিকোনো}}, \emph{bikono}.) Goods? (\textbf{\bn{মাল}}, \emph{māl}.) A warehouse? (\textbf{\bn{গুদাম}}, \emph{gudām}.) An order (for goods or food)? (\textbf{\bn{অর্ডার}}, \emph{ôrḍār}.) Wholesale? (\textbf{\bn{পাইকারি}}, \emph{pāikāri}.) To stock up? (\textbf{\bn{মজুত} \bn{করা}}, \emph{môjut kôrā}.) A pair of scales? (\textbf{\bn{দাঁড়িপাল্লা}}, \emph{dā̃ṛipāllā}.) Pure (unmixed)? (\textbf{\bn{খাঁটি}}, \emph{khā̃ṭi}.) A shopping list? (\textbf{\bn{ফর্দ}}, \emph{phôrdô}.) A customer (formal)? (\textbf{\bn{গ্রাহক}}, \emph{grāhôk}.) To test (by trying)? (\textbf{\bn{পরখ} \bn{করা}}, \emph{pôrôkh kôrā}.) A paper bag? (\textbf{\bn{ঠোঙা}}, \emph{ṭhongā}.) A wrapper? (\textbf{\bn{মোড়ক}}, \emph{moṛôk}.) A carrier bag? (\textbf{\bn{ক্যারিব্যাগ}}, \emph{kæribyāg}.) A size (of clothes or shoes)? (\textbf{\bn{সাইজ}}, \emph{sāij}.) To palm off? (\textbf{\bn{গছানো}}, \emph{gôchhāno}.)
\end{tcolorbox}
